\chapter{Processus de développement : modèle en Y}

\section{Principe}

Le projet suit le \textbf{modèle en Y} : un processus itératif, piloté par les cas d'utilisation, dans lequel la
capture des besoins fonctionnels et la capture des besoins techniques sont menées en parallèle avant de se
rejoindre dans la branche de conception et de réalisation (figure~\ref{fig:modele-y}).

\begin{figure}[H]
  \centering
  \begin{tikzpicture}[
      etape/.style={draw=primaire,fill=fondgris,rounded corners=3pt,thick,
                    align=center,text width=6.2cm,inner sep=7pt,font=\small},
      fleche/.style={-{Stealth[length=3mm]},thick,draw=secondaire}]
    \node[etape] (fonc) at (-3.8,0)
      {\textbf{Branche fonctionnelle}\\besoins, cas d'utilisation,\\classes du domaine};
    \node[etape] (tech) at (3.8,0)
      {\textbf{Branche technique}\\Java 17, Jakarta EE 10, Tomcat 10.1,\\Hibernate 6, MySQL 8, architecture MVC};
    \node[etape] (conc) at (0,-2.7)
      {\textbf{Conception préliminaire et détaillée}\\classes de conception, séquences};
    \node[etape] (real) at (0,-5.0)
      {\textbf{Codage, tests, déploiement (.war)}\\itérations successives,\\une branche par fonctionnalité};
    \draw[fleche] (fonc.south) -- (conc.north west);
    \draw[fleche] (tech.south) -- (conc.north east);
    \draw[fleche] (conc) -- (real);
  \end{tikzpicture}
  \caption{Le modèle en Y appliqué au projet}\label{fig:modele-y}
\end{figure}

\section{Branche fonctionnelle}

\begin{itemize}
  \item Rédaction du cahier des charges (\code{docs/cahier-des-charges.md}) : 23 besoins fonctionnels codés BF01 à BF23
        et priorisés (P1 indispensable, P2 important, P3 souhaitable), 16 règles de gestion RG01 à RG16,
        12 besoins non fonctionnels BNF01 à BNF12.
  \item Diagrammes de cas d'utilisation (Visiteur/Client et Administrateur), descriptions textuelles
        (chapitre~\ref{chap:cas}), diagramme de classes du domaine.
\end{itemize}

\section{Branche technique}

\begin{itemize}
  \item Technologies imposées par le module : Java 17, Jakarta EE 10 (Servlet 6.0, JSP 3.1, JSTL 3.0), Tomcat 10.1,
        Hibernate ORM 6, MySQL 8, Maven (packaging \code{war}), Bootstrap 5, jBCrypt, PlantUML.
  \item Architecture MVC N-tiers : navigateur, Tomcat (couches web et métier), serveur MySQL.
  \item Configuration Hibernate (\code{hibernate.cfg.xml}, génération des tables par \code{hbm2ddl.auto = update}) et
        configuration web (\code{web.xml} : session, encodage, pages d'erreur).
\end{itemize}

\section{Branche de conception et de réalisation}

\begin{itemize}
  \item Conception : diagrammes de classes de conception (couche web, couche métier) et diagrammes de séquence.
        Tous les diagrammes ont été versionnés avant la première ligne de code (commit \code{4fc0fac}), puis mis à jour
        au fil des itérations quand l'implémentation a précisé la conception (commits \code{docs:} des branches
        \code{feature/catalogue}, \code{feature/commande}, \code{feature/gestion-commandes}, \code{feature/p3} et
        \code{feature/coherence-finale}).
  \item Réalisation : une itération par groupe de fonctionnalités, chacune livrée sur une branche \code{feature/*},
        vérifiée par des tests manuels (navigateur et requêtes HTTP) puis fusionnée dans \code{main} par pull request.
\end{itemize}

\section{Itérations et suivi GitHub}

Le backlog est suivi dans \textbf{GitHub Issues} : un ticket par étape, avec la liste à cocher des codes BF/RG
concernés. Chaque étape est développée sur sa propre branche et fusionnée par une pull request dont la
description reprend les besoins couverts, les règles appliquées et les tests manuels effectués
(tableau~\ref{tab:iterations}).

\begin{table}[H]
  \centering\small
  \caption{Itérations (cahier des charges §8.3), tickets, branches et pull requests}\label{tab:iterations}
  \begin{tabularx}{\linewidth}{cLcP{4.5cm}c}
    \ligneentete
    \entete{It.} & \entete{Étape} & \entete{Ticket} & \entete{Branche} & \entete{PR} \\
    1 & Modélisation UML (cas d'utilisation, classes, séquences) & \#4 & \code{main} (documentation) & — \\
    1 & Squelette : Maven, Hibernate, entités, encodage & \#5 & \code{feature/squelette} & \#1 \\
    2 & Authentification et filtres (BF01 à BF03) & \#6 & \code{feature/authentification} & \#2 \\
    2 & Back-office catégories, produits, stock (BF17 à BF19) & \#7 & \code{feature/admin-catalogue} & \#3 \\
    2 & Catalogue public (BF05 à BF08) & \#8 & \code{feature/catalogue} & \#14 \\
    3 & Panier en session (BF09 à BF12) & \#9 & \code{feature/panier} & \#15 \\
    3 & Passage de commande (BF13, RG09 à RG12) & \#10 & \code{feature/commande} & \#16 \\
    3 & Historique et gestion des commandes (BF14, BF15, BF20, BF21) & \#11 & \code{feature/gestion-commandes} & \#17 \\
    4 & Utilisateurs (BF22), pages d'erreur, validation, responsive & \#12 & \code{feature/finitions} & \#18 \\
    4 & Besoins P3 : profil (BF04), annulation client (BF16), tableau de bord (BF23) & \#12 & \code{feature/p3} & \#19 \\
    5 & README et rapport & \#13 & \code{feature/livrables} & \#20 \\
    5 & Cohérence finale : séquences, ordre des filtres, documentation & \#21 & \code{feature/coherence-finale} & \#22 \\
    \bottomrule
  \end{tabularx}
\end{table}

Remarques :
\begin{itemize}
  \item Les tickets et les pull requests partagent la même numérotation sur GitHub, d'où les numéros non
        consécutifs des PR (\#1 à \#3, puis \#14 à \#20, puis \#22).
  \item Le texte « Closes \#N » en fin de description de PR ferme automatiquement le ticket à la fusion.
  \item Les P3 ont été regroupés dans une PR séparée (\#19) pour pouvoir être isolés de l'étape 9 si nécessaire.
\end{itemize}

\section{Conventions Git}

\begin{itemize}
  \item Messages de commit en français, sous la forme \code{type: description} : \code{feat} (fonctionnalité), \code{fix}
        (correction), \code{docs} (documentation, UML), \code{style} (présentation), \code{chore} (configuration du projet).
        Exemples : \code{feat: passage de commande transactionnel (RG09 à RG12)},
        \code{fix: réintégration du stock calculée par la base à l'annulation (RG14)}.
  \item Un commit par sujet ; les références BF/RG figurent dans le message quand c'est utile.
  \item Fusion avec un commit de fusion explicite : \code{fusion: <description> (\#N)}, puis suppression de la branche.
\end{itemize}
